\documentclass{article}
\usepackage{amsmath, amssymb, amsthm}
\usepackage{hyperref}
\usepackage{geometry}
\geometry{margin=1in}

\title{Erd\H{o}s Problem \#782}
\date{Last updated: 2025-08-31}
\author{Source: erdosproblems.com}

\begin{document}
\maketitle

\noindent\textbf{Status:} OPEN \quad \textbf{No prize}

\noindent\textbf{Tags:} number theory

\medskip
\noindent\textbf{Problem Statement:}

Do the squares contain arbitrarily long quasi-progressions? That is, does there exist some constant $C>0$ such that, for any $k$, the squares contain a sequence $x_1,\ldots,x_k$ where, for some $d$ and all $1\leq i<k$,\[x_i+d\leq x_{i+1}\leq x_i+d+C.\]Do the squares contain arbitrarily large cubes\[a+\left\{ \sum_i \epsilon_ib_i : \epsilon_i\in \{0,1\}\right\}?\]

\bigskip
\noindent\textbf{Remarks:}

A question of Brown, Erd\H{o}s, and Freedman \cite{BEF90}. It is a classical fact that the squares do not contain arithmetic progressions of length $4$.

An affirmative answer to the first question implies an affirmative answer to the second.

Solymosi \cite{So07} conjectured the answer to the second question is no. Cilleruelo and Granville \cite{CiGr07} have observed that the answer to the second question is no conditional on the Bombieri-Lang conjecture.

\bigskip
\noindent\textbf{References:}

[BEF90] Brown, T. C. and Erd\H{o}s, P. and Freedman, A. R., Quasi-progressions and descending waves. J. Combin. Theory Ser. A (1990), 81-95.
 
 [CiGr07] Cilleruelo, Javier and Granville, Andrew, Lattice points on circles, squares in arithmetic progressions
and sumsets of squares. (2007), 241-262.
 
 [So07] Solymosi, J\'{o}zsef, Elementary additive combinatorics. (2007), 29-38.

\bigskip
\noindent\small{Source: \url{https://www.erdosproblems.com/782}}
\end{document}
